% =====================================================================
%  chapters/minggu-04.tex
%  Minggu 4 — HTML Lanjutan: Form, Tabel & Multimedia
% =====================================================================

\chapter{HTML Lanjutan: Form, Tabel \& Multimedia}
\label{chap:minggu-04}

% ---------------------------------------------------------------------
% 1. CAPAIAN PEMBELAJARAN
% ---------------------------------------------------------------------
\begin{capaian}
Setelah menyelesaikan bab ini, mahasiswa diharapkan mampu:
\begin{itemize}
  \item membuat form HTML dengan berbagai jenis \emph{input} (\emph{text}, \emph{email}, \emph{password}, \emph{checkbox}, \emph{radio}, \emph{select}, \emph{textarea});
  \item menggunakan \texttt{<label>}, \texttt{<fieldset>}, dan validasi form dasar;
  \item membuat tabel HTML yang terstruktur dan aksesibel (\texttt{<thead>}, \texttt{<tbody>}, \texttt{scope});
  \item menyisipkan elemen multimedia (gambar, video, audio) ke halaman HTML;
  \item menggabungkan semua elemen ke dalam satu halaman yang rapi.
\end{itemize}
\end{capaian}

% ---------------------------------------------------------------------
% 2. TUJUAN PRAKTIK
% ---------------------------------------------------------------------
\begin{tujuanpraktik}
Pada sesi praktik ini, mahasiswa akan:
\begin{itemize}
  \item membangun form pendaftaran akun dengan berbagai jenis \emph{input};
  \item membuat tabel data mahasiswa dengan struktur aksesibel;
  \item menyisipkan gambar, video, dan audio ke halaman HTML;
  \item menggabungkan semua elemen ke dalam satu halaman praktikum.
\end{itemize}
\end{tujuanpraktik}

% ---------------------------------------------------------------------
% 3. FORM HTML
% ---------------------------------------------------------------------
\section{Form HTML}
\label{sec:m04-form}

\begin{catatan}
\textbf{Form} adalah cara pengguna mengirim data ke server. Kita fokus pada struktur form-nya dulu; pengiriman data (\emph{submit}) akan dipelajari saat belajar JavaScript.
\end{catatan}

\subsection{Live Coding: Form Pendaftaran Akun}

\begin{langkah}
  \item Di folder \texttt{latihan-web-3} (atau buat folder baru \texttt{latihan-web-4}), buat file \texttt{form.html}.
  \item Ketik struktur HTML dasar + hubungkan ke \texttt{style.css} (dari Bab~\ref{chap:minggu-03}), atau buat style \emph{inline} baru.
  \item Ketik kode form berikut di dalam \texttt{<body>}.
\end{langkah}

\begin{lstlisting}[language=HTML, caption={Form pendaftaran akun}, label={lst:m04-form}]
<!DOCTYPE html>
<html lang="id">
<head>
    <meta charset="UTF-8">
    <meta name="viewport"
          content="width=device-width, initial-scale=1.0">
    <title>Form Pendaftaran --- Latihan HTML Lanjutan</title>
    <style>
        body {
            font-family: Arial, sans-serif;
            max-width: 600px;
            margin: 0 auto;
            padding: 20px;
            line-height: 1.6;
        }
        h1 { color: #2c3e50; }
        form {
            background: #f9f9f9;
            padding: 20px;
            border-radius: 8px;
        }
        label {
            display: block;
            margin-top: 12px;
            font-weight: bold;
        }
        input, select, textarea {
            width: 100%;
            padding: 8px;
            margin-top: 4px;
            border: 1px solid #ccc;
            border-radius: 4px;
            box-sizing: border-box;
        }
        fieldset {
            border: 1px solid #ccc;
            border-radius: 8px;
            padding: 12px;
            margin-top: 16px;
        }
        legend {
            font-weight: bold;
            padding: 0 8px;
        }
        .radio-group label,
        .checkbox-group label {
            display: inline;
            font-weight: normal;
        }
        button {
            margin-top: 16px;
            padding: 10px 20px;
            background: #2c3e50;
            color: white;
            border: none;
            border-radius: 4px;
            font-size: 16px;
            cursor: pointer;
        }
        button:hover { background: #34495e; }
    </style>
</head>
<body>

    <h1>Form Pendaftaran Akun</h1>
    <p>Isi form berikut untuk mendaftar.
       Tanda <span style="color:red;">*</span>
       wajib diisi.</p>

    <form action="#" method="post">

        <!-- FIELDSET 1: Data Diri -->
        <fieldset>
            <legend>Data Diri</legend>

            <label for="nama">Nama Lengkap
                <span style="color:red;">*</span></label>
            <input type="text" id="nama" name="nama"
                   placeholder="Masukkan nama lengkap"
                   required>

            <label for="email">Email
                <span style="color:red;">*</span></label>
            <input type="email" id="email" name="email"
                   placeholder="contoh@email.com"
                   required>

            <label for="password">Password
                <span style="color:red;">*</span></label>
            <input type="password" id="password"
                   name="password"
                   placeholder="Minimal 8 karakter"
                   minlength="8" required>

            <label for="tanggal">Tanggal Lahir</label>
            <input type="date" id="tanggal" name="tanggal">

            <label for="jk">Jenis Kelamin
                <span style="color:red;">*</span></label>
            <div class="radio-group">
                <input type="radio" id="laki" name="jk"
                       value="laki-laki" required>
                <label for="laki">Laki-laki</label>

                <input type="radio" id="perempuan"
                       name="jk" value="perempuan">
                <label for="perempuan">Perempuan</label>
            </div>
        </fieldset>

        <!-- FIELDSET 2: Informasi Akademik -->
        <fieldset>
            <legend>Informasi Akademik</legend>

            <label for="prodi">Program Studi
                <span style="color:red;">*</span></label>
            <select id="prodi" name="prodi" required>
                <option value="">-- Pilih Prodi --</option>
                <option value="d3-informatika">
                    D3 Informatika</option>
                <option value="d3-multimedia">
                    D3 Multimedia</option>
                <option value="d3-akuntansi">
                    D3 Akuntansi</option>
                <option value="d3-administrasi">
                    D3 Administrasi Bisnis</option>
            </select>

            <label for="angkatan">Angkatan</label>
            <input type="number" id="angkatan"
                   name="angkatan" min="2020" max="2030"
                   placeholder="Contoh: 2026">

            <label>Hobi / Minat (boleh pilih lebih dari
                   1):</label>
            <div class="checkbox-group">
                <input type="checkbox" id="hobi-web"
                       name="hobi" value="web-design">
                <label for="hobi-web">Desain Web</label>
                <br>
                <input type="checkbox" id="hobi-mobile"
                       name="hobi" value="mobile-app">
                <label for="hobi-mobile">Mobile App</label>
                <br>
                <input type="checkbox" id="hobi-ai"
                       name="hobi" value="ai">
                <label for="hobi-ai">AI / Machine
                    Learning</label>
                <br>
                <input type="checkbox" id="hobi-network"
                       name="hobi" value="networking">
                <label for="hobi-network">Jaringan</label>
            </div>
        </fieldset>

        <!-- FIELDSET 3: Pesan -->
        <fieldset>
            <legend>Catatan Tambahan</legend>

            <label for="pesan">Tulis pesan / motivasi
                   kamu:</label>
            <textarea id="pesan" name="pesan" rows="5"
                      placeholder="Tulis di sini...">
            </textarea>
        </fieldset>

        <!-- Tombol Submit -->
        <button type="submit">Daftar Sekarang</button>

    </form>

</body>
</html>
\end{lstlisting}

Simpan $\rightarrow$ buka dengan \emph{Live Server}.

\subsection{Jenis-Jenis Input}

\begin{table}[ht]
\centering
\caption{Jenis input form HTML}
\label{tab:m04-input}
\small
\begin{tabular}{lp{3.5cm}p{5cm}}
\toprule
\textbf{Jenis Input} & \textbf{Fungsi} & \textbf{Contoh Penggunaan} \\
\midrule
\texttt{type="text"} & Teks biasa (satu baris) & Nama, judul \\
\texttt{type="email"} & Teks + validasi format email & Alamat email \\
\texttt{type="password"} & Teks tersembunyi (bullet) & Password \\
\texttt{type="number"} & Hanya angka + tombol naik/turun & Usia, tahun \\
\texttt{type="date"} & \emph{Date picker} (kalender) & Tanggal lahir \\
\texttt{type="radio"} & Pilih satu dari beberapa opsi & Jenis kelamin \\
\texttt{type="checkbox"} & Pilih beberapa opsi & Hobi, minat \\
\texttt{type="tel"} & Nomor telepon & No.\ HP \\
\texttt{type="url"} & Teks + validasi format URL & Website \\
\texttt{type="file"} & \emph{Upload} file & Upload foto profil \\
\texttt{<select>} & \emph{Dropdown} (daftar pilihan) & Program studi \\
\texttt{<textarea>} & Teks panjang (multi-baris) & Pesan, komentar \\
\bottomrule
\end{tabular}
\end{table}

\subsection{Pentingnya Label}

\begin{lstlisting}[language=HTML, caption={Perbandingan label benar dan salah}, label={lst:m04-label}]
<!-- Benar: label terhubung ke input via "for" = id -->
<label for="nama">Nama Lengkap</label>
<input type="text" id="nama" name="nama">

<!-- Salah: label tidak terhubung -->
<label>Nama Lengkap</label>
<input type="text">
\end{lstlisting}

\begin{catatan}
Kenapa harus pakai label?
\begin{enumerate}
  \item \textbf{Aksesibilitas:} \emph{Screen reader} bisa membaca ``Nama Lengkap'' saat input difokuskan.
  \item \textbf{UX:} Klik pada label akan memfokuskan kursor ke input yang sesuai.
  \item \textbf{Validasi:} Peramban menampilkan pesan error yang lebih informatif.
\end{enumerate}
\end{catatan}

\subsection{Atribut Validasi Dasar}

\begin{table}[ht]
\centering
\caption{Atribut validasi form}
\label{tab:m04-validasi}
\begin{tabular}{lp{4cm}p{5cm}}
\toprule
\textbf{Atribut} & \textbf{Fungsi} & \textbf{Contoh} \\
\midrule
\texttt{required} & Input wajib diisi & Semua field penting \\
\texttt{minlength} / \texttt{maxlength} & Batas panjang teks & Password minimal 8 karakter \\
\texttt{min} / \texttt{max} & Batas angka & Tahun 2020--2030 \\
\texttt{placeholder} & Teks petunjuk di dalam input & ``Masukkan email'' \\
\texttt{pattern} & Pola \emph{regex} untuk validasi & Format telepon tertentu \\
\texttt{disabled} & Nonaktifkan input & Field yang belum tersedia \\
\texttt{readonly} & Boleh dilihat, tidak bisa diedit & Field otomatis \\
\bottomrule
\end{tabular}
\end{table}

\subsection{Verifikasi Form}

Coba:
\begin{enumerate}
  \item Klik tombol \textbf{Daftar Sekarang} tanpa mengisi apapun $\rightarrow$ peramban harus menampilkan pesan error ``Please fill out this field'' pada field \texttt{required}.
  \item Isi \emph{email} dengan format salah (misal: \texttt{bukanemail}) $\rightarrow$ tekan \emph{Submit} $\rightarrow$ peramban harus menampilkan pesan error tentang format email.
  \item Isi \emph{password} kurang dari 8 karakter $\rightarrow$ peramban harus menolak.
\end{enumerate}

\begin{tangkapanlayar}
  {assets/images/minggu-04/form-pendaftaran.png}
  {Form pendaftaran di peramban.}
  {Pastikan terlihat semua field (Data Diri, Info Akademik, Pesan) dan tombol ``Daftar Sekarang''.}
\end{tangkapanlayar}

% ---------------------------------------------------------------------
% 4. TABEL HTML
% ---------------------------------------------------------------------
\section{Tabel HTML}
\label{sec:m04-tabel}

\begin{catatan}
\textbf{Tabel} digunakan untuk menampilkan data berbentuk baris dan kolom. Penting untuk data terstruktur seperti jadwal, perbandingan harga, statistik, dsb.
\end{catatan}

\subsection{Live Coding: Tabel Sederhana}

Buat file baru: \texttt{tabel.html}. Ketik kode berikut:

\begin{lstlisting}[language=HTML, caption={Tabel data mahasiswa dengan colspan dan rowspan}, label={lst:m04-tabel}]
<!DOCTYPE html>
<html lang="id">
<head>
    <meta charset="UTF-8">
    <meta name="viewport"
          content="width=device-width, initial-scale=1.0">
    <title>Tabel Data --- Latihan HTML Lanjutan</title>
    <style>
        body {
            font-family: Arial, sans-serif;
            max-width: 800px;
            margin: 20px auto;
        }
        table {
            width: 100%;
            border-collapse: collapse;
            margin: 20px 0;
        }
        th, td {
            border: 1px solid #ddd;
            padding: 10px 12px;
            text-align: left;
        }
        thead { background-color: #2c3e50; color: white; }
        tbody tr:nth-child(even) {
            background-color: #f2f2f2;
        }
        caption {
            font-size: 1.2em;
            font-weight: bold;
            margin-bottom: 10px;
        }
    </style>
</head>
<body>

    <h1>Data Mahasiswa D3 Informatika</h1>

    <!-- TABLE 1: Tabel Dasar -->
    <table>
        <caption>Jadwal Kuliah Semester 1</caption>
        <thead>
            <tr>
                <th scope="col">No</th>
                <th scope="col">Mata Kuliah</th>
                <th scope="col">Hari</th>
                <th scope="col">Jam</th>
                <th scope="col">SKS</th>
            </tr>
        </thead>
        <tbody>
            <tr>
                <td>1</td>
                <td>Desain Web</td>
                <td>Senin</td>
                <td>08:00 &ndash; 10:00</td>
                <td>3</td>
            </tr>
            <tr>
                <td>2</td>
                <td>Algoritma Pemrograman</td>
                <td>Selasa</td>
                <td>10:00 &ndash; 12:00</td>
                <td>3</td>
            </tr>
            <tr>
                <td>3</td>
                <td>Basis Data</td>
                <td>Rabu</td>
                <td>13:00 &ndash; 15:00</td>
                <td>3</td>
            </tr>
            <tr>
                <td>4</td>
                <td>Matematika Diskrit</td>
                <td>Kamis</td>
                <td>08:00 &ndash; 10:00</td>
                <td>3</td>
            </tr>
            <tr>
                <td>5</td>
                <td>Bahasa Inggris Teknik</td>
                <td>Jumat</td>
                <td>10:00 &ndash; 12:00</td>
                <td>2</td>
            </tr>
        </tbody>
        <tfoot>
            <tr>
                <td colspan="4"><strong>Total SKS</strong></td>
                <td><strong>14</strong></td>
            </tr>
        </tfoot>
    </table>

    <hr>

    <!-- TABLE 2: Tabel dengan colspan & rowspan -->
    <h2>Perbandingan Nilai UTS</h2>
    <table>
        <thead>
            <tr>
                <th rowspan="2" scope="col">Mata Kuliah</th>
                <th colspan="3"
                    scope="colgroup">Nilai Per Kelompok</th>
            </tr>
            <tr>
                <th scope="col">Kelompok A</th>
                <th scope="col">Kelompok B</th>
                <th scope="col">Kelompok C</th>
            </tr>
        </thead>
        <tbody>
            <tr>
                <td>Desain Web</td>
                <td>85</td>
                <td>78</td>
                <td>90</td>
            </tr>
            <tr>
                <td>Algoritma</td>
                <td>72</td>
                <td>88</td>
                <td>81</td>
            </tr>
            <tr>
                <td>Basis Data</td>
                <td>91</td>
                <td>76</td>
                <td>85</td>
            </tr>
        </tbody>
    </table>

</body>
</html>
\end{lstlisting}

Simpan $\rightarrow$ buka dengan \emph{Live Server}.

\subsection{Struktur Tabel}

\begin{lstlisting}[caption={Struktur dasar tabel HTML}, label={lst:m04-struktur-tabel}]
<table>
    <caption>   Judul tabel (wajib ada untuk aksesibilitas)
    <thead>     Bagian kepala tabel (baris judul kolom)
        <tr>    Table row (baris)
        <th>    Table header (judul kolom/baris)
    <tbody>     Bagian isi tabel (data)
        <tr>
        <td>    Table data (sel data)
    <tfoot>     Bagian kaki tabel (total, ringkasan)
        <tr>
        <td>
</table>
\end{lstlisting}

\subsection{Atribut Penting untuk Aksesibilitas}

\begin{table}[ht]
\centering
\caption{Atribut tabel untuk aksesibilitas}
\label{tab:m04-tabel-attr}
\begin{tabular}{lp{3.5cm}p{5.5cm}}
\toprule
\textbf{Atribut} & \textbf{Fungsi} & \textbf{Contoh} \\
\midrule
\texttt{scope="col"} & Menandai header sebagai kolom & \texttt{<th scope="col">Nama</th>} \\
\texttt{scope="row"} & Menandai header sebagai baris & \texttt{<th scope="row">Senin</th>} \\
\texttt{colspan="N"} & Sel melebar N kolom & \texttt{<td colspan="3">Gabungan</td>} \\
\texttt{rowspan="N"} & Sel melebar N baris & \texttt{<th rowspan="2">Header</th>} \\
\bottomrule
\end{tabular}
\end{table}

\begin{catatan}
Atribut \texttt{scope} membantu \emph{screen reader} menavigasi tabel dengan benar.
\end{catatan}

\begin{tangkapanlayar}
  {assets/images/minggu-04/tabel-data.png}
  {Tabel data mahasiswa di peramban.}
  {Pastikan terlihat 2 tabel: Jadwal Kuliah dan Perbandingan Nilai, dengan \emph{styling} yang rapi.}
\end{tangkapanlayar}

% ---------------------------------------------------------------------
% 5. MULTIMEDIA
% ---------------------------------------------------------------------
\section{Multimedia: Gambar, Video, Audio}
\label{sec:m04-multimedia}

\subsection{Gambar --- Review \& Lanjutan}

Kita sudah belajar \texttt{<img>} di Bab~\ref{chap:minggu-03}. Sekarang tambahkan \texttt{<figure>}:

\begin{lstlisting}[language=HTML, caption={Gambar dengan figure dan figcaption}, label={lst:m04-figure}]
<figure>
    <img src="https://placehold.co/600x400"
         alt="Pemandangan gunung saat matahari terbit"
         width="600" height="400"
         loading="lazy">
    <figcaption>Gambar 1 --- Pemandangan alam sebagai
        contoh konten visual</figcaption>
</figure>
\end{lstlisting}

Atribut tambahan \texttt{loading="lazy"}: gambar hanya dimuat saat pengguna \emph{scroll} ke bagian tersebut (hemat \emph{bandwidth} \& mempercepat \emph{loading} halaman pertama).

\subsection{Video (HTML5)}

\begin{lstlisting}[language=HTML, caption={Menyisipkan video HTML5}, label={lst:m04-video}]
<!-- Video dari URL langsung -->
<video width="640" height="360" controls>
    <source src="video-tutorial.mp4" type="video/mp4">
    Browser kamu tidak mendukung elemen video.
</video>

<!-- Contoh dengan video publik -->
<video width="640" height="360" controls
       poster="https://placehold.co/640x360">
    <source src="https://www.w3schools.com/html/mov_bbb.mp4"
            type="video/mp4">
    Browser kamu tidak mendukung elemen video.
</video>
\end{lstlisting}

\begin{table}[ht]
\centering
\caption{Atribut elemen video}
\label{tab:m04-video}
\begin{tabular}{lp{8cm}}
\toprule
\textbf{Atribut} & \textbf{Fungsi} \\
\midrule
\texttt{controls} & Menampilkan tombol play/pause/volume \\
\texttt{autoplay} & Video otomatis dimuat (hindari --- mengganggu pengguna) \\
\texttt{loop} & Video diulang terus-menerus \\
\texttt{muted} & Video tanpa suara \\
\texttt{poster} & Gambar \emph{thumbnail} sebelum video dimainkan \\
\texttt{width} / \texttt{height} & Ukuran tampilan video \\
\bottomrule
\end{tabular}
\end{table}

\subsection{Audio (HTML5)}

\begin{lstlisting}[language=HTML, caption={Menyisipkan elemen audio}, label={lst:m04-audio}]
<audio controls>
    <source src="https://www.w3schools.com/html/horse.mp3"
            type="audio/mpeg">
    Browser kamu tidak mendukung elemen audio.
</audio>
\end{lstlisting}

\begin{table}[ht]
\centering
\caption{Atribut elemen audio}
\label{tab:m04-audio}
\begin{tabular}{lp{8cm}}
\toprule
\textbf{Atribut} & \textbf{Fungsi} \\
\midrule
\texttt{controls} & Menampilkan tombol play/pause/volume \\
\texttt{autoplay} & Audio otomatis diputar (sangat tidak disarankan) \\
\texttt{loop} & Audio diulang terus-menerus \\
\texttt{muted} & Audio tanpa suara \\
\bottomrule
\end{tabular}
\end{table}

\subsection{Iframe --- \emph{Embed} Konten Eksternal}

\begin{lstlisting}[language=HTML, caption={Menyisipkan video YouTube via iframe}, label={lst:m04-iframe}]
<!-- Embed YouTube video -->
<iframe width="560" height="315"
    src="https://www.youtube.com/embed/dQw4w9WgXcQ"
    title="Contoh video YouTube"
    frameborder="0"
    allow="accelerometer; autoplay; clipboard-write;
           encrypted-media; gyroscope; picture-in-picture"
    allowfullscreen>
</iframe>
\end{lstlisting}

\begin{peringatan}
Selalu isi atribut \texttt{title} pada \texttt{<iframe>} untuk aksesibilitas.
\end{peringatan}

% ---------------------------------------------------------------------
% 6. PRAKTIK: MENGGabungkan SEMUA ELEMEN
% ---------------------------------------------------------------------
\section{Praktik: Menggabungkan Semua Elemen}
\label{sec:m04-gabung}

\begin{langkah}
  \item Buat file \texttt{praktikum-minggu-4.html} yang menggabungkan semua elemen yang sudah dipelajari.
  \item Gunakan struktur \emph{semantic} HTML (\texttt{<header>}, \texttt{<nav>}, \texttt{<main>}, \texttt{<section>}, \texttt{<footer>}).
  \item Sertakan section form, tabel, dan multimedia dalam satu halaman.
\end{langkah}

\begin{lstlisting}[language=HTML, caption={Halaman praktikum minggu 4 (ringkasan struktur)}, label={lst:m04-gabung}]
<!DOCTYPE html>
<html lang="id">
<head>
    <meta charset="UTF-8">
    <meta name="viewport"
          content="width=device-width, initial-scale=1.0">
    <title>Praktikum Minggu 4 --- Form, Tabel &amp;
           Multimedia</title>
    <style>
        * { box-sizing: border-box; }
        body {
            font-family: Arial, sans-serif;
            max-width: 900px;
            margin: 0 auto;
            padding: 20px;
            line-height: 1.6;
            color: #333;
        }
        h1 {
            color: #2c3e50;
            border-bottom: 2px solid #2c3e50;
            padding-bottom: 8px;
        }
        h2 { color: #34495e; margin-top: 40px; }
        nav {
            background: #2c3e50;
            padding: 12px;
            border-radius: 8px;
            margin-bottom: 20px;
        }
        nav a {
            color: white;
            text-decoration: none;
            margin: 0 12px;
        }
        nav a:hover { text-decoration: underline; }
        table {
            width: 100%;
            border-collapse: collapse;
            margin: 16px 0;
        }
        th, td {
            border: 1px solid #ddd;
            padding: 8px 12px;
            text-align: left;
        }
        thead { background: #2c3e50; color: white; }
        caption {
            font-weight: bold;
            margin-bottom: 8px;
        }
        form {
            background: #f9f9f9;
            padding: 20px;
            border-radius: 8px;
        }
        label {
            display: block;
            margin-top: 10px;
            font-weight: bold;
        }
        input, select, textarea {
            width: 100%;
            padding: 8px;
            margin-top: 4px;
            border: 1px solid #ccc;
            border-radius: 4px;
        }
        button {
            margin-top: 14px;
            padding: 10px 24px;
            background: #2c3e50;
            color: white;
            border: none;
            border-radius: 4px;
            cursor: pointer;
            font-size: 16px;
        }
        button:hover { background: #34495e; }
        figure {
            margin: 20px 0;
            text-align: center;
        }
        figcaption {
            font-style: italic;
            margin-top: 6px;
            color: #666;
        }
        footer {
            background: #2c3e50;
            color: white;
            text-align: center;
            padding: 12px;
            border-radius: 8px;
            margin-top: 40px;
        }
    </style>
</head>
<body>

    <!-- HEADER & NAV -->
    <header>
        <h1>Praktikum Minggu 4</h1>
        <nav>
            <a href="#form">Form</a>
            <a href="#tabel">Tabel</a>
            <a href="#multimedia">Multimedia</a>
        </nav>
    </header>

    <main>

        <!-- SECTION: FORM -->
        <section id="form">
            <h2>1. Form Pendaftaran</h2>
            <form action="#" method="post">
                <label for="p-nama">Nama Lengkap
                    <span style="color:red;">*</span></label>
                <input type="text" id="p-nama" name="nama"
                       placeholder="Nama kamu" required>

                <label for="p-email">Email
                    <span style="color:red;">*</span></label>
                <input type="email" id="p-email"
                       name="email"
                       placeholder="email@contoh.com"
                       required>

                <label for="p-prodi">Program Studi</label>
                <select id="p-prodi" name="prodi">
                    <option value="">-- Pilih --</option>
                    <option value="d3-informatika">
                        D3 Informatika</option>
                    <option value="d3-multimedia">
                        D3 Multimedia</option>
                </select>

                <label for="p-pesan">Pesan</label>
                <textarea id="p-pesan" name="pesan"
                          rows="4"
                          placeholder="Tulis pesan...">
                </textarea>

                <button type="submit">Kirim</button>
            </form>
        </section>

        <!-- SECTION: TABEL -->
        <section id="tabel">
            <h2>2. Tabel Perbandingan</h2>
            <table>
                <caption>Perbandingan Framework CSS</caption>
                <thead>
                    <tr>
                        <th scope="col">Fitur</th>
                        <th scope="col">CSS Biasa</th>
                        <th scope="col">Tailwind CSS</th>
                    </tr>
                </thead>
                <tbody>
                    <tr>
                        <td>Cara penulisan</td>
                        <td>File terpisah (.css)</td>
                        <td>Class inline di HTML</td>
                    </tr>
                    <tr>
                        <td>Learning curve</td>
                        <td>Rendah</td>
                        <td>Sedang</td>
                    </tr>
                    <tr>
                        <td>Ukuran file akhir</td>
                        <td>Terkontrol manual</td>
                        <td>Di-browse otomatis</td>
                    </tr>
                </tbody>
            </table>
        </section>

        <!-- SECTION: MULTIMEDIA -->
        <section id="multimedia">
            <h2>3. Multimedia</h2>

            <figure>
                <img src="https://placehold.co/600x300"
                     alt="Contoh gambar placeholder"
                     width="600" height="300"
                     loading="lazy">
                <figcaption>Gambar 1 --- Demonstrasi
                    elemen &lt;img&gt; dengan
                    &lt;figure&gt;</figcaption>
            </figure>

            <h3>Video</h3>
            <video width="100%" controls
                   poster="https://placehold.co/640x360">
                <source
                    src="https://www.w3schools.com/html/mov_bbb.mp4"
                    type="video/mp4">
                Browser kamu tidak mendukung elemen video.
            </video>

            <h3>Audio</h3>
            <audio controls style="width:100%;">
                <source
                    src="https://www.w3schools.com/html/horse.mp3"
                    type="audio/mpeg">
                Browser kamu tidak mendukung elemen audio.
            </audio>
        </section>

    </main>

    <footer>
        <p>&copy; 2026 --- Praktikum Minggu 4 |
           D3 Informatika</p>
    </footer>

</body>
</html>
\end{lstlisting}

\begin{tip}
Pastikan di peramban:
\begin{itemize}
  \item Tautan \emph{anchor} (\#form, \#tabel, \#multimedia) berfungsi.
  \item Form bisa diisi dan menampilkan pesan error jika field \emph{required} kosong.
  \item Tabel muncul dengan benar.
  \item Gambar muncul dengan \emph{caption}.
  \item Video dan audio bisa diputar.
  \item \emph{Footer} muncul di bagian bawah.
\end{itemize}
\end{tip}

\begin{tangkapanlayar}
  {assets/images/minggu-04/praktikum-section-form.png}
  {Section Form di halaman praktikum.}
  {Pastikan terlihat label, input field, select \emph{dropdown}, textarea, dan tombol \emph{submit}.}
\end{tangkapanlayar}

\begin{tangkapanlayar}
  {assets/images/minggu-04/praktikum-section-tabel-multimedia.png}
  {Section Tabel dan Multimedia.}
  {Pastikan terlihat tabel perbandingan, gambar dengan \emph{caption}, video \emph{player}, dan audio \emph{player}.}
\end{tangkapanlayar}

% ---------------------------------------------------------------------
% 7. VALIDASI & PUSH KE GITHUB
% ---------------------------------------------------------------------
\section{Validasi \& Push ke GitHub}
\label{sec:m04-push}

\subsection{Validasi HTML}

\begin{langkah}
  \item Buka \texttt{https://validator.w3.org/\#validate\_by\_input}.
  \item Salin-tempel isi \texttt{praktikum-minggu-4.html} $\rightarrow$ klik \textbf{Check}.
  \item Perbaiki semua error hingga hasilnya bersih.
\end{langkah}

\subsection{Push ke GitHub}

\begin{lstlisting}[caption={Perintah push ke GitHub}, label={lst:m04-git}]
# Inisialisasi (jika folder baru)
git init
git remote add origin https://github.com/NAMA_USER/latihan-web-4.git

# Commit dan push
git add .
git commit -m "Latihan HTML minggu 4: form, tabel, multimedia"
git branch -M main
git push -u origin main
\end{lstlisting}

% ---------------------------------------------------------------------
% 8. LATIHAN MANDIRI
% ---------------------------------------------------------------------
\section{Latihan Mandiri}
\label{sec:m04-latihan}

\begin{latihan}[Form ``Survei Kepuasan'']
Buat halaman \texttt{survei.html} yang berisi form survei kepuasan mahasiswa terhadap layanan kampus. Field yang wajib ada:
\begin{enumerate}
  \item \textbf{Nama} --- \texttt{type="text"}, \texttt{required}
  \item \textbf{NIM} --- \texttt{type="number"}, \texttt{required}, \texttt{minlength="8"}
  \item \textbf{Email} --- \texttt{type="email"}, \texttt{required}
  \item \textbf{Fakultas} --- \texttt{<select>} (minimal 4 pilihan), \texttt{required}
  \item \textbf{Semester saat ini} --- \texttt{type="number"}, \texttt{min="1"} \texttt{max="8"}
  \item \textbf{Tingkat kepuasan} --- \texttt{type="radio"} (1--5), \texttt{required}, gunakan \texttt{<fieldset>} + \texttt{<legend>}
  \item \textbf{Saran/keluhan} --- \texttt{<textarea>}, \texttt{rows="5"}
  \item \textbf{Setuju dihubungi} --- \texttt{type="checkbox"}
\end{enumerate}
Gunakan minimal 2 \texttt{<fieldset>} untuk mengelompokkan field. Setiap field harus punya \texttt{<label>} yang terhubung (\texttt{for} = \texttt{id}). Gunakan \texttt{<main>}, \texttt{<header>}, \texttt{<footer>} (\emph{semantic HTML}).
\end{latihan}

\begin{latihan}[Tabel ``Data Buku'']
Buat halaman \texttt{buku.html} yang berisi tabel daftar buku dengan kolom: No, Judul Buku, Pengarang, Tahun, Harga. Persyaratan:
\begin{itemize}
  \item Minimal 8 baris data (buku).
  \item Gunakan \texttt{<thead>}, \texttt{<tbody>}, \texttt{<tfoot>} (\texttt{<tfoot>} untuk total rata-rata harga).
  \item Gunakan \texttt{scope="col"} pada header kolom.
  \item Gunakan \texttt{colspan} di \emph{footer} untuk kolom ``Total''.
  \item Tambahkan \texttt{<caption>} pada tabel.
  \item Sertakan minimal 1 gambar (misal: sampul buku \emph{placeholder}) di atas tabel dengan \texttt{<figure>} + \texttt{<figcaption>}.
\end{itemize}
\end{latihan}

\begin{latihan}[Push \& Screenshot]
\begin{enumerate}
  \item \emph{Push} semua file ke GitHub.
  \item Ambil \emph{ screenshot} hasil latihan (form \& tabel).
  \item Pastikan tidak ada \emph{warning} di W3C Validator.
\end{enumerate}
\end{latihan}

% ---------------------------------------------------------------------
% 9. TROUBLESHOOTING
% ---------------------------------------------------------------------
\section{Troubleshooting}
\label{sec:m04-troubleshoot}

\begin{table}[ht]
\centering
\caption{Masalah umum form, tabel, dan multimedia}
\label{tab:m04-troubleshoot}
\small
\begin{tabular}{p{4.5cm}p{8cm}}
\toprule
\textbf{Masalah} & \textbf{Solusi} \\
\midrule
Form tidak melakukan apa-apa saat \emph{submit} & \texttt{action="\#"} hanya \emph{placeholder} --- tidak ada \emph{backend}. Validasi peramban tetap berjalan meskipun tidak ada server \\
\emph{Radio button} bisa dipilih semua & Pastikan semua \emph{radio button} punya nama yang sama (\texttt{name="jk"}), sehingga hanya satu yang bisa dipilih \\
Label tidak memfokuskan input saat diklik & Pastikan atribut \texttt{for} pada \texttt{<label>} = \texttt{id} pada \texttt{<input>} (harus sama persis) \\
Tampilan tabel berantakan di \emph{mobile} & Tabel memang tidak responsif secara \emph{default} --- akan dipelajari cara membuatnya responsif di CSS \\
Video/audio tidak muncul & Periksa URL/path file. Jika dari URL eksternal, pastikan URL masih aktif \\
\texttt{<caption>} tidak terlihat & Pastikan \texttt{<caption>} adalah anak langsung \texttt{<table>} dan berada sebelum \texttt{<thead>} \\
Tombol \emph{submit} me-\emph{reload} halaman & Normal --- karena \texttt{action="\#"}; saat belajar JavaScript, \emph{reload} akan dicegah dengan \texttt{event.preventDefault()} \\
\bottomrule
\end{tabular}
\end{table}

% ---------------------------------------------------------------------
% 10. RANGKUMAN
% ---------------------------------------------------------------------
\section{Rangkuman}
\label{sec:m04-rangkuman}

\begin{itemize}
  \item Membuat form HTML dengan berbagai jenis \emph{input}.
  \item Menggunakan \texttt{<label>}, \texttt{<fieldset>}, \texttt{<legend>} untuk form yang rapi \& aksesibel.
  \item Atribut validasi: \texttt{required}, \texttt{minlength}, \texttt{min}, \texttt{max}, \texttt{placeholder}.
  \item Membuat tabel dengan \texttt{<table>}, \texttt{<caption>}, \texttt{<thead>}, \texttt{<tbody>}, \texttt{<tfoot>}.
  \item Atribut tabel: \texttt{scope}, \texttt{colspan}, \texttt{rowspan}.
  \item Menyisipkan multimedia: \texttt{<img>}, \texttt{<figure>}, \texttt{<video>}, \texttt{<audio>}, \texttt{<iframe>}.
\end{itemize}

\begin{tip}
\textbf{Cheat Sheet Form Input Types:}
\begin{itemize}
  \item \texttt{text} $\rightarrow$ Teks biasa
  \item \texttt{email} $\rightarrow$ Email (validasi format)
  \item \texttt{password} $\rightarrow$ Teks tersembunyi
  \item \texttt{number} $\rightarrow$ Angka + tombol +/--
  \item \texttt{date} $\rightarrow$ \emph{Date picker} (kalender)
  \item \texttt{tel} $\rightarrow$ Nomor telepon
  \item \texttt{url} $\rightarrow$ URL (validasi format)
  \item \texttt{radio} $\rightarrow$ Pilih 1 dari beberapa
  \item \texttt{checkbox} $\rightarrow$ Pilih beberapa
  \item \texttt{file} $\rightarrow$ \emph{Upload} file
  \item \texttt{color} $\rightarrow$ \emph{Color picker}
  \item \texttt{range} $\rightarrow$ Slider
  \item \texttt{hidden} $\rightarrow$ Tersembunyi
  \item \texttt{submit} $\rightarrow$ Tombol \emph{submit}
  \item \texttt{reset} $\rightarrow$ Tombol \emph{reset} form
\end{itemize}
\end{tip}
